%% sections/05_problem_formulation.tex

\section{예측 조건: 언제 spec-decode가 이득인가}
\label{sec:problem}

\subsection{방법-이득과 지배 방법의 부재}

방법 $m \in \mathcal{M}$ (vanilla, suffix, ngram, llm-d)의
(모델 $\theta$, 워크로드 $w$)에서의 처리량 이득을 정의한다.

\begin{definition}[방법-이득]
\label{def:method-gain}
\[
\Delta(m, w, \theta) = \frac{T(m, w, \theta) - T(\mathrm{vanilla}, w, \theta)}
                            {T(\mathrm{vanilla}, w, \theta)},
\]
$\Delta(\mathrm{vanilla},\cdot)=0$ 이고 $\Delta<0$ 이면 $m$ 이 net-negative 다.
\end{definition}

\begin{property}[지배 방법의 부재]
\label{prop:no-dominant}
모든 $(w,\theta)$ 에 대해 $\Delta(m^*,w,\theta)>0$ 인 단일
$m^*\in\mathcal{M}\setminus\{\mathrm{vanilla}\}$ 는 존재하지 않는다.
\end{property}

\noindent\textbf{근거(B200, \S\ref{sec:results}):}
suffix 는 70셀 중 55셀에서 net-positive 이나 15셀에서 vanilla 에 미달하고
(추론$\cdot$MoE 모델 전반), naive n-gram 은 실 트레이스에서 $\sim\!-90\%$ 로 붕괴하며,
llm-d 는 최선 단일 방법을 14\%(10/70)의 셀에서만 능가한다.
어떤 단일 방법도 보편 최적이 아니다.

\subsection{부호의 예측 조건: $\alpha$ 대 overhead}

식~\ref{eq:rk-ratio}로부터 spec-decode 가 이득인 필요충분 조건은 다음이다.

\begin{equation}
\Delta(\mathrm{suffix}) > 0 \iff \alpha\,K_{\mathrm{eff}} > \mathrm{overhead}(\theta).
\label{eq:rk-condition}
\end{equation}

두 항이 부호를 가른다.
\textbf{(i) 수락률 $\alpha$} — 드래프트 품질$\times$출력 예측 가능성.
B200 실측에서 net-positive 셀의 $\alpha$ 중앙값은 $0.68$,
vanilla-win 셀은 $0.37$ 이다(\S\ref{subsec:res-alpha}).
\textbf{(ii) overhead$(\theta)$} — 모델 규모와 아키텍처(MoE expert 활성)에 비례.
따라서 break-even $\alpha$ 는 \emph{고정값이 아니라 모델에 따라 이동}한다:
overhead 가 작은 모델은 낮은 $\alpha$ 에서도 이기고(예: 한 셀은 $\alpha\!=\!0.28$ 에서 $+209\%$),
overhead 가 큰 R1-671B(MoE)는 $\alpha\!\approx\!0.45\text{--}0.54$ 에서도 $-40\!\sim\!-64\%$ 손해다.

\smallskip
\noindent\textbf{문제 A (모델-유형 인지 방법 선택).}
회귀($\Delta<0$)를 vanilla 로 차단하면서 가속 가능 시 $\Delta>0$ 방법을 택하는
경량 게이트 $g:(\theta,w)\to\mathcal{M}$ 를 설계한다.
\S\ref{sec:results} 는 부호의 지배 축이 워크로드가 아니라 \emph{모델 유형}
(추론/MoE 여부)임을 보이며, 이는 $g$ 의 1차 입력이 모델 프로파일임을 함의한다.

\subsection{문제 B: 성공이 만드는 host 병목}
\label{subsec:problem-b}

spec-decode 가 성공하면 step 수가 줄어 GPU 가 un-saturate 된다.

\begin{definition}[host 병목 전환]
\label{def:host-bound}
$\Delta(\mathrm{suffix})>0$ 인 셀에서 GPU 이용률이 vanilla 대비 하락하면
($u_{\mathrm{gpu}}^{\mathrm{suf}} < u_{\mathrm{gpu}}^{\mathrm{van}}$),
줄어든 step 사이의 host 작업(드래프팅$\cdot$스케줄$\cdot$샘플링$\cdot$detok)이
새로운 병목 후보가 된다. 회수 가능 slack 의 절대량은
$u_{\mathrm{cpu}}\ll1$ 인 유휴 CPU 와 $1-u_{\mathrm{gpu}}^{\mathrm{suf}}$ 의 GPU 여유다.
\end{definition}

B200 실측에서 net-positive 셀의 GPU 이용률은 $94.9\%\!\to\!78.7\%$(중앙)로 떨어지고
CPU 는 전 구성에서 $\le4.9\%$ 유휴다(\S\ref{subsec:res-resource}).

\smallskip
\noindent\textbf{문제 B (host-side slack reclamation).}
spec-decode 가 만든 host-path 병목을 유휴 CPU 로 떠안아
GPU 여유를 다시 처리량으로 환산한다.
이는 CPU 로 GPU 연산을 흉내내는 것이 아니라(\S\ref{sec:direction}),
GPU 가 \emph{내려놓은} host 작업을 CPU 가 받는 것이다.
